\chapter*{Conclusion Générale}
\addcontentsline{toc}{chapter}{Conclusion Générale}

Le projet \textbf{VoteUASZ} répond au besoin de dématérialisation des scrutins de l'Université Assane Seck de Ziguinchor en proposant une application web complète, de l'import CSV des électeurs à la génération du procès-verbal PDF.

Sur le plan technique, l'architecture sépare l'émargement (\texttt{VoterAuditLog}) du bulletin (\texttt{Ballot} sans clé étrangère vers \texttt{User}), chiffre les bulletins en \textbf{AES-256-GCM} (PBKDF2), relie la chaîne des bulletins en \textbf{SHA-256}, sécurise l'accès par \textbf{JWT HS256} avec \textbf{Refresh Token} rotatif et OTP e-mail, et diffuse les résultats en temps réel via \textbf{STOMP} et \textbf{Chart.js}. La qualité est vérifiée par \textbf{122 tests backend} (couverture JaCoCo \textbf{74,30\,\%} en instructions, \textbf{0 violation Checkstyle}) et \textbf{26 tests frontend} (\textbf{0 échec}).

De manière lucide, certaines garanties restent perfectibles (clé AES portée par le serveur, corrélation d'ordre possible entre l'ID séquentiel d'émargement et \texttt{sequenceNumber}, absence de signature RSA/X.509 du PV PDF), ouvrant la voie à des perspectives claires pour une version ultérieure.

% ============================================================================
% BIBLIOGRAPHIE
% ============================================================================
\begin{thebibliography}{9}
\addcontentsline{toc}{chapter}{Bibliographie}

\bibitem{spring-boot}
VMware Tanzu, \textit{Spring Boot 3.2.3 Reference Documentation}, 2024. \url{https://docs.spring.io/spring-boot/docs/3.2.3/reference/html/}

\bibitem{spring-security}
VMware Tanzu, \textit{Spring Security Reference}, 2024. \url{https://docs.spring.io/spring-security/reference/}

\bibitem{angular}
Google, \textit{Angular 17 Documentation}, 2024. \url{https://angular.io/docs}

\bibitem{postgresql}
PostgreSQL Global Development Group, \textit{PostgreSQL 16 Documentation}, 2024. \url{https://www.postgresql.org/docs/16/}

\bibitem{nist-gcm}
NIST, \textit{Recommendation for Block Cipher Modes of Operation: Galois/Counter Mode (GCM) -- SP 800-38D}, 2007.

\bibitem{rfc7519}
M. Jones, J. Bradley, N. Sakimura, \textit{JSON Web Token (JWT) -- RFC 7519}, IETF, 2015.

\bibitem{openpdf}
LibrePDF, \textit{OpenPDF 1.3.30 Documentation}, GitHub. \url{https://github.com/LibrePDF/OpenPDF}

\bibitem{owasp}
OWASP Foundation, \textit{OWASP Top 10 Web Application Security Risks}, 2021. \url{https://owasp.org/www-project-top-ten/}

\end{thebibliography}

% ============================================================================
% ANNEXES
% ============================================================================
\appendix

\chapter{Localisation des algorithmes et contrat OpenAPI}

\begin{table}[H]
\centering
\caption{Localisation des algorithmes clés dans le backend}
\small
\begin{tabularx}{\textwidth}{@{}lXl@{}}
\toprule
\rowcolor{uaszlight}
\textbf{Mécanisme} & \textbf{Fichier} & \textbf{Lignes} \\
\midrule
Vérification OTP \& \texttt{VoteToken} & \texttt{service/VotingService.java} & \texttt{L94--L142} \\
Transaction de vote \& chaînage SHA-256 & \texttt{service/VotingService.java} & \texttt{L150--L228} \\
Chiffrement AES-256-GCM \& PBKDF2 & \texttt{service/CryptoService.java} & \texttt{L62--L82, L130--L139} \\
Vérification chaîne de hash & \texttt{service/AuditService.java} & \texttt{L52--L77} \\
Refresh token \& mot de passe & \texttt{controller/AuthController.java} & \texttt{L107--L240} \\
Machine à états (6 états) & \texttt{service/ElectionStateMachine.java} & \texttt{L10--L35} \\
\bottomrule
\end{tabularx}
\end{table}

\begin{table}[H]
\centering
\caption{Principaux endpoints REST (\texttt{docs/openapi.yaml}, préfixe \texttt{/api/v1})}
\small
\begin{tabularx}{\textwidth}{@{}llX@{}}
\toprule
\rowcolor{uaszlight}
\textbf{Module} & \textbf{Endpoints} & \textbf{Accès} \\
\midrule
Authentification & \texttt{POST /auth/login}, \texttt{/refresh}, \texttt{/logout}, \texttt{/forgot-password}, \texttt{/reset-password} & Public (sauf \texttt{change-password}) \\
Scrutins & \texttt{GET /elections}, \texttt{POST/PUT/DELETE /elections}, \texttt{PUT /elections/\{id\}/status} & Public en lecture ; \texttt{SUPER\_ADMIN} en écriture \\
Candidatures & \texttt{POST /candidacies}, \texttt{GET /candidacies/election/\{id\}}, \texttt{PUT /.../validate} & \texttt{CANDIDAT} ; \texttt{COMMISSION} / \texttt{ADMIN} \\
Campagnes & \texttt{POST /campaigns}, \texttt{GET /campaigns/election/\{id\}} & \texttt{CANDIDAT} approuvé ; authentifié \\
Vote OTP & \texttt{POST /vote/request-otp}, \texttt{/verify-otp}, \texttt{/submit}, \texttt{GET /vote/status/\{id\}} & Électeur éligible authentifié \\
Résultats \& PV & \texttt{GET /results/\{id\}}, \texttt{GET /results/\{id\}/export-pdf} & Selon statut ; \texttt{ADMIN} / \texttt{COMMISSION} \\
Réclamations & \texttt{POST /complaints}, \texttt{GET /complaints/*}, \texttt{PUT /complaints/\{id\}/resolve} & Authentifié ; \texttt{COMMISSION} / \texttt{ADMIN} \\
Audit \& Import & \texttt{GET /audit/\{id\}}, \texttt{/export-csv}, \texttt{POST /admin/users/import-csv} & \texttt{SUPER\_ADMIN} / \texttt{COMMISSION} \\
\bottomrule
\end{tabularx}
\end{table}

\chapter{Guide d'installation locale et exemple d'audit / PV}

\section{Démarrage en local (Docker Compose + Angular)}
\begin{enumerate}
    \item Copier \texttt{.env.example} vers \texttt{.env} et renseigner \texttt{POSTGRES\_PASSWORD}, \texttt{JWT\_SECRET} (32+ car.), \texttt{CRYPTO\_SECRET} (32+ car.), \texttt{DEFAULT\_ELECTEUR\_PASSWORD} et \texttt{DEFAULT\_CANDIDAT\_PASSWORD}.
    \item Démarrer PostgreSQL 16 (port 5540) et MailHog (SMTP 1025, UI 8025) : \texttt{docker compose up -d postgres mailhog}.
    \item Lancer le backend Java 21 : \texttt{mvn spring-boot:run} (ou \texttt{java -jar vote-uasz-backend-1.0.0.jar}) dans \texttt{backend/}.
    \item Lancer le frontend Angular 17 : \texttt{cd frontend \&\& npm ci \&\& npx ng serve --port 4200}.
\end{enumerate}

\section{Exemple d'export CSV d'audit et de PV PDF}
Lors du jeu de données local de démonstration (scrutin \textit{Élection du délégué de L3 Informatique 2026}, 25 votes enregistrés), l'appel à \texttt{AuditService.verifyBallotChain()} retourne \texttt{chainValid = true} avec 25 bulletins chaînés depuis \texttt{"GENESIS"} et 25 émargements anonymisés. Le procès-verbal généré pour \textit{Élection du Vice-recteur 2026} (\texttt{CLOTURE}) est présenté en figure~\ref{fig:pv-pdf} (chapitre 3).
